\documentclass[10pt,a4paper]{amsart}
\usepackage[margin=19mm]{geometry}
\usepackage{amsmath,amssymb,mathtools}
\usepackage[scr=rsfs,cal=euler]{mathalfa}
\usepackage{xfrac}
\usepackage[unicode,hidelinks,bookmarks=false]{hyperref}
\newcommand{\ie}{i.e.\ }
\newcommand{\cf}{cf.\ }
\newcommand{\resp}{respectively\ }
\newcommand{\Subsection}{\subsection}
\providecommand{\nobreakdash}{\nobreak\hbox{-}}
\setlength{\emergencystretch}{2em}
\multlinegap0pt
\usepackage{bm}
\usepackage{bbm}
\usepackage{dsfont}
\usepackage{xspace}
\usepackage{cancel}
\usepackage{mathtools}
\usepackage{amsthm}
\makeatletter
\@ifundefined{proposition}{\newtheorem{proposition}{Proposition}}{}
\@ifundefined{remark}{\newtheorem{remark}{Remark}}{}
\makeatother
\newcommand{\Nn}{{\mathcal N}}
\newcommand{\g}{{\, |\,}}
\newcommand{\oT}{\boldsymbol{T}}
\newcommand{\wt}{{\widetilde t}}
\title[Random tree Besov priors: Data-driven regularisation paramet]{Random tree Besov priors: Data-driven regularisation parameter selection}
\author{Hanne Kekkonen, Andreas Tataris}
\date{Source: arXiv:2601.12957v1 (CC BY 4.0)}
\begin{document}
\maketitle

\noindent\emph{Source and licence.} Random tree Besov priors: Data-driven regularisation parameter selection. Version arXiv:2601.12957v1 (CC BY 4.0), distributed under the Creative Commons Attribution 4.0 International licence, as recorded in the arXiv licence metadata for this exact version and shown on the abstract page https://arxiv.org/abs/2601.12957v1. Full rights notice: licenses/arxiv-2601.12957-CC-BY-4.0.txt.\par\smallskip

\noindent\textit{Editorial scope.} Counted unit: the Remark whose source label is \texttt{None} in the article (source0.tex lines 595--600, source label \texttt{None}), delivered together with the article's own complete original proof of it (source0.tex lines 602--642, one \texttt{proof} environment of 41 lines). No mathematics is written, repaired or re-derived: statement and proof are the article's own text line for line. Required context retained uncounted: eq:DenoisingMain (lines 290--292); eq:DefPrior (lines 327--332); Prop:scaling (lines 579--587). Every definition, display and label of those lines is retained unchanged. Omitted from this package and not redistributed: 1--289, 293--326, 333--578, 588--594, 601--601, 643--815. Nothing inside a retained excerpt is omitted or altered: every retained body is delivered exactly as the article's lines are written, so no omission receipt of any kind accompanies this package. Citations: the retained excerpts carry no citation command at all, so no bibliography entry of the article is transcribed and no reference is redistributed. Reference adaptation: none was needed and none was made; every reference inside the delivered unit resolves inside it, so no unresolved reference or citation is produced. Printed slips kept literal: no printed word, symbol, hypothesis or number of the counted statement or of its proof was altered. Layout and class: the article's own class is \texttt{article} with packages including abstract, amsmath, amssymb, amsthm, hyperref, booktabs, caption, geometry, subcaption, graphicx, pgfplots, adjustbox, algorithm2e, nowidow; the wrapper is the lane's pinned \texttt{amsart} editorial wrapper at 10pt on A4, which restates the theorem-like environments the retained text needs and adds the article's own macro definitions that the retained text needs (\texttt{\textbackslash Nn}, \texttt{\textbackslash g}, \texttt{\textbackslash oT}, \texttt{\textbackslash wt}), with extra wrapper packages (bm, bbm, dsfont, xspace, cancel, mathtools). The delivered numbering is the unit's own. Assets: no retained span contains a figure environment, a graphics-inclusion command, a PDF-page inclusion, a table or a TikZ picture; no image file is copied into this package, the package declares no asset dependency and no image file is redistributed with it. Licence: version arXiv:2601.12957v1 of this article is distributed under the Creative Commons Attribution 4.0 International licence, as recorded in the arXiv licence metadata for this exact version and shown on the abstract page https://arxiv.org/abs/2601.12957v1. A full rights notice is delivered at licenses/arxiv-2601.12957-CC-BY-4.0.txt. Characters: every retained excerpt is pure ASCII, so the delivered engine, XeLaTeX with its default Latin Modern text font, reports no missing character.
\begin{align}\label{eq:DenoisingMain}
    M_i=f(x_i)+W_i,
\end{align}

\begin{align}\label{eq:DefPrior}
\begin{split}
f(x) & =\sum_{(j,k)\in T} \langle f,\psi_{jk}\rangle\psi_{jk}(x)\\
& =\sum_{(j,k)\in \oT} \wt_{jk} g_{jk}\psi_{jk}(x),
\end{split}
\end{align}

\begin{proposition}\label{Prop:scaling}
Consider the denoising problem of estimating $f$ from a noisy measurement \eqref{eq:DenoisingMain}. We assume a truncated random tree Besov prior for $f$ as described in \eqref{eq:DefPrior}. Denote by $T^{1}_\beta(m_1)$ the denoised signal arising from tree pruning with a fixed $\beta$ when $w_i\sim \Nn(0,1)$ and $g_{jk}\sim\Nn(0,1)$. 

Then, for $\sigma,\kappa>0$, the denoised signal $T^{\kappa}_\beta(m_\sigma)$ arising from a measurement with noise $w_i\sim\Nn(0,\sigma^2)$ and prior choice $g_{jk}\sim\Nn(0,\kappa^2)$ is given by $\tilde{T}^1_{\tilde{\beta}}(m_1)$, where the coefficients $\hat{g}^1_{jk}=m_{jk}/2$ are replaced by $\hat{g}^1_{jk}=\kappa^2m_{jk}/(\kappa^2+\sigma^2)$ and
\begin{align*}
    \widetilde{\beta}=\frac{\beta^{c}}{\beta^{c}+(1-\beta)^{c}}
\end{align*}
with $c=\sigma^2(\kappa^2+\sigma^2)/(2\kappa^2)>0$.
\end{proposition}

\begin{remark}
We see from Proposition \ref{Prop:scaling} that, for a constant $c$, $T^1_\beta(cm_1) =cT^1_{\widetilde{\beta}}(m_1)$ where 
\begin{align*}
    \widetilde{\beta}=\frac{\beta^{c^2}}{\beta^{c^2}+(1-\beta)^{c^2}}.
\end{align*}
\end{remark}

\begin{proof}[Proof of Proposition \ref{Prop:scaling}]
We start by noting that when $\beta$ is fixed and $w_{jk}\sim\Nn(0,\sigma^2)$, $g_{jk}\sim\Nn(0,\kappa^2)$ the posterior of $g_{jk}\g \tilde{t}_{jk}=1$ is given by
\begin{align*}
z_{jk}^1 & \propto \exp\Big(-\frac{1}{2\sigma^2}(m_{jk}-g_{jk})^2
-\frac{1}{2\kappa^2}g_{jk}^2)\Big)\beta
\end{align*}
and the posterior for $g_{jk}\g \tilde{t}_{jk}=0$ is 
\begin{align*}
z_{jk}^0 & \propto \exp\Big(-\frac{1}{2\sigma^2}m_{jk}^2-\frac{1}{2\kappa^2}g_{jk}^2)\Big)(1-\beta).
\end{align*}

In the bottom row $j=J$ the weight resulting from including a node $(J,k)$ is simply
\begin{align*}
    Y_{Jk}^1=\min_{g_{Jk}}\Big(-\log(z_{Jk}^1)\Big)
=\frac{m_{Jk}^2}{2(\kappa^2+\sigma^2)}-\log(\beta),
\end{align*}
which is achieved at $\hat{g}_{Jk}=\kappa^2m_{Jk}/(\kappa^2+\sigma^2)$, and the weight for not including the node is 
\begin{align*}
    Y_{Jk}^0=\min_{g_{Jk}}\Big(-\log(z_{Jk}^0)\Big)
=\frac{m_{Jk}^2}{2\sigma^2}-\log(1-\beta),
\end{align*}
following from the choice $\hat{g}_{Jk}=0$.

A node $(J,k)$ is included in the subtree if including the node carries less weight than turning it off, that is, when $Y_{Jk}^1\leq Y_{Jk}^0$ or equivalently when 
\begin{align*}
    \log\Big(\frac{1-\beta}{\beta}\Big)
    \leq \frac{\kappa^2m_{Jk}^2}{2\sigma^2(\kappa^2+\sigma^2)}
    =\frac{1}{c}\frac{m_{Jk}^2}{4},
\end{align*}
where $c=\sigma^2(\kappa^2+\sigma^2)/(2\kappa^2)>0$.  We can rewrite the above as 
\begin{align*}
    \log\Bigg(\Big(\frac{1-\beta}{\beta}\Big)^c\Bigg)
    =\frac{m_{Jk}^2}{4},
\end{align*}
which corresponds to 
\begin{align*}
    \log\Big(\frac{1-\widetilde{\beta}}{\widetilde{\beta}}\Big)
    =\frac{m_{Jk}^2}{4},
\end{align*}
where $\widetilde{\beta}=\beta^c/(\beta^c+(1-\beta)^c)\in[0,0.5]$ for $\beta\in[0,0.5]$ and $c>0$. 
\end{proof}

\end{document}
